\documentclass[12pt, a4paper]{book}

% ─── Packages ───────────────────────────────────────────────────────────────
\usepackage[a4paper, margin=2.5cm]{geometry}
\usepackage{fontenc}
\usepackage{inputenc}
\usepackage{microtype}
\usepackage{graphicx}
\usepackage{xcolor}
\usepackage{titlesec}
\usepackage{titletoc}
\usepackage{enumitem}
\usepackage{booktabs}
\usepackage{array}
\usepackage{tabularx}
\usepackage{fancyhdr}
\usepackage{hyperref}
\usepackage{mdframed}
\usepackage{multicol}
\usepackage{parskip}

% ─── Colours ────────────────────────────────────────────────────────────────
\definecolor{primary}{HTML}{2E7D32}      % deep green
\definecolor{accent}{HTML}{FF8F00}       % warm amber
\definecolor{lightgray}{HTML}{F5F5F5}
\definecolor{darkgray}{HTML}{424242}

% ─── Hyperref ───────────────────────────────────────────────────────────────
\hypersetup{
    colorlinks=true,
    linkcolor=primary,
    urlcolor=accent
}

% ─── Header / Footer ────────────────────────────────────────────────────────
\pagestyle{fancy}
\fancyhf{}
\fancyhead[LE,RO]{\small\textcolor{darkgray}{\leftmark}}
\fancyhead[RE,LO]{\small\textcolor{primary}{\textbf{Jorge's Recipe Book}}}
\fancyfoot[C]{\textcolor{darkgray}{\thepage}}
\renewcommand{\headrulewidth}{0.4pt}

% ─── Section / Chapter Styling ──────────────────────────────────────────────
\titleformat{\chapter}[display]
  {\normalfont\huge\bfseries\color{primary}}
  {}
  {0pt}
  {\Huge}
  [{\vspace{0.5ex}\titlerule[1.5pt]}]

\titleformat{\section}
  {\normalfont\Large\bfseries\color{primary}}
  {}
  {0pt}
  {}
  [{\vspace{0.2ex}\textcolor{accent}{\rule{\linewidth}{1pt}}}]

% ─── Custom Environments ─────────────────────────────────────────────────────

% Macro box (coloured info box)
\newmdenv[
  backgroundcolor=lightgray,
  linecolor=primary,
  linewidth=1.5pt,
  roundcorner=6pt,
  skipabove=10pt,
  skipbelow=10pt,
  innerleftmargin=12pt,
  innerrightmargin=12pt,
  innertopmargin=10pt,
  innerbottommargin=10pt
]{macrobox}

% Tip box
\newmdenv[
  backgroundcolor=accent!10,
  linecolor=accent,
  linewidth=1.5pt,
  roundcorner=6pt,
  skipabove=10pt,
  skipbelow=10pt,
  innerleftmargin=12pt,
  innerrightmargin=12pt,
  innertopmargin=10pt,
  innerbottommargin=10pt
]{tipbox}

% ─── Helper Commands ─────────────────────────────────────────────────────────
\newcommand{\recipetitle}[1]{%
  \section*{\textcolor{primary}{#1}}%
  \addcontentsline{toc}{section}{#1}%
}

\newcommand{\recipemeta}[3]{%
  \begin{macrobox}
    \begin{tabular}{lll}
      \textbf{\textcolor{primary}{Cost per serve:}} & #1 &
      \textbf{\textcolor{primary}{Calories:}} #2 \quad
      \textbf{\textcolor{primary}{Protein:}} #3
    \end{tabular}
  \end{macrobox}%
}

\newcommand{\ingredientheader}{\subsection*{\textcolor{accent}{Ingredients}}}
\newcommand{\methodheader}{\subsection*{\textcolor{accent}{Method}}}
\newcommand{\noteheader}{\subsection*{\textcolor{accent}{Notes}}}

% ════════════════════════════════════════════════════════════════════════════
\begin{document}

% ─── Title Page ──────────────────────────────────────────────────────────────
\begin{titlepage}
  \centering
  \vspace*{3cm}
  {\fontsize{48}{56}\selectfont\bfseries\textcolor{primary}{Jorge's}\par}
  \vspace{0.3cm}
  {\fontsize{32}{40}\selectfont\bfseries\textcolor{accent}{Recipe Book}\par}
  \vspace{0.5cm}
  {\large\textcolor{darkgray}{A 4-Week Meal Plan for a Healthier You}\par}
  \vspace{1.5cm}
  \textcolor{primary}{\rule{0.6\linewidth}{2pt}}
  \vspace{1.5cm}

  {\large\textcolor{darkgray}{
    39 recipes designed around your lifestyle ---\\[6pt]
    high protein, budget friendly, and easy to batch cook.
  }\par}

  \vfill
  {\small\textcolor{darkgray}{Christchurch, New Zealand \\ \the\year}}
\end{titlepage}

% ─── Table of Contents ───────────────────────────────────────────────────────
\tableofcontents
\clearpage

% ════════════════════════════════════════════════════════════════════════════
%  4-WEEK PLAN OVERVIEW
% ════════════════════════════════════════════════════════════════════════════
\chapter*{Your 4-Week Plan}
\addcontentsline{toc}{chapter}{Your 4-Week Plan}

\begin{tipbox}
  \textbf{How your week works:}\\
  \textbf{Monday \& Thursday} --- Flatmate cooks for you.\\
  \textbf{Tuesday \& Friday} --- You cook for both of you (6 serves each).\\
  \textbf{Wednesday \& Saturday} --- Leftovers from the night before.\\
  \textbf{Sunday} --- Youth group dinner. No cooking needed!
\end{tipbox}

\vspace{0.5cm}

\begin{tabularx}{\linewidth}{|l|X|X|}
  \hline
  \rowcolor{primary!20}
  \textbf{Week} & \textbf{Tuesday} & \textbf{Friday} \\
  \hline
  Week 1 & Slow Cooker Mexican Chicken & Beef Mince \& Lentil Shepherd's Pie \\
  \hline
  Week 2 & Beef Mince Stir Fry & Slow Cooker Chicken \& Chickpea Stew \\
  \hline
  Week 3 & Beef \& Vegetable Chilli & Slow Cooker Pulled Pork Wraps \\
  \hline
  Week 4 & Slow Cooker Honey Garlic Chicken & Pork \& Lentil Bolognese \\
  \hline
  \multicolumn{3}{|l|}{\textit{Bonus meals to rotate in: Beef \& Broccoli Stir Fry, Slow Cooker Beef \& Vegetable Soup}} \\
  \hline
\end{tabularx}

\vspace{1cm}

\subsection*{\textcolor{primary}{Breakfast Rotation}}

\begin{tabularx}{\linewidth}{|l|X|}
  \hline
  \rowcolor{primary!20}
  \textbf{Week} & \textbf{Breakfast Prep} \\
  \hline
  Week 1 & Egg Muffins \\
  \hline
  Week 2 & Overnight Oats \\
  \hline
  Week 3 & Breakfast Burritos \\
  \hline
  Week 4 & Frittata Slices \\
  \hline
\end{tabularx}

\vspace{1cm}

\begin{macrobox}
  \textbf{\textcolor{primary}{Daily Targets}} \\[4pt]
  \begin{tabular}{ll}
    Calories: & 2,000--2,200 kcal/day \\
    Protein:  & 150g/day \\
    Steps:    & 8,000/day \\
  \end{tabular}\\[6pt]
  \textit{Note: Two meals alone won't quite reach these targets. Add a protein shake, Greek yogurt, or a handful of nuts as a daily snack.}
\end{macrobox}

% ════════════════════════════════════════════════════════════════════════════
%  DINNER RECIPES
% ════════════════════════════════════════════════════════════════════════════
\chapter{Dinner Recipes}

% ─── Recipe 1 ────────────────────────────────────────────────────────────────
\recipetitle{1. Slow Cooker Mexican Chicken}
\recipemeta{\$2.20/serve}{470 kcal}{48g}

\ingredientheader
\begin{itemize}[leftmargin=1.5em]
  \item 1 kg chicken thighs
  \item 1 can black beans, drained and rinsed
  \item 1 can corn, drained
  \item 1 jar medium salsa (400g)
  \item 1 tsp cumin
  \item 1 tsp smoked paprika
  \item 1 tsp garlic powder
  \item \textbf{To serve:} rice, Greek yogurt, shredded cheese
\end{itemize}

\methodheader
\begin{enumerate}[leftmargin=1.5em]
  \item Place chicken thighs in the slow cooker.
  \item Add black beans, corn, salsa, and all spices. Stir to combine.
  \item Cook on \textbf{low for 6--8 hours} or high for 3--4 hours.
  \item Shred the chicken with two forks directly in the slow cooker and stir through the sauce.
  \item Serve over rice with a dollop of Greek yogurt and shredded cheese.
\end{enumerate}

\begin{tipbox}
  \textbf{Tip:} Makes 6 serves. Leftovers keep in the fridge for 3--4 days or freeze for up to 3 months.
\end{tipbox}

\clearpage

% ─── Recipe 2 ────────────────────────────────────────────────────────────────
\recipetitle{2. Beef Mince \& Lentil Shepherd's Pie}
\recipemeta{\$2.80/serve}{480 kcal}{42g}

\ingredientheader
\begin{itemize}[leftmargin=1.5em]
  \item 500g lean beef mince
  \item 1 cup red lentils, rinsed
  \item 2 cans diced tomatoes
  \item 1 onion, diced
  \item 2 carrots, diced
  \item 2 stalks celery, diced
  \item 1 tbsp Worcestershire sauce
  \item 1 cup beef stock
  \item 4 large kumara (sweet potato), peeled and cubed
  \item Salt, pepper, butter for mash
\end{itemize}

\methodheader
\begin{enumerate}[leftmargin=1.5em]
  \item Preheat oven to 180\textdegree C.
  \item Brown mince in a large oven-safe pan over medium-high heat. Drain excess fat.
  \item Add onion, carrots, and celery. Cook for 5 minutes until softened.
  \item Stir in lentils, diced tomatoes, Worcestershire sauce, and beef stock. Simmer for 20 minutes until lentils are soft.
  \item Meanwhile, boil kumara until tender, then mash with butter, salt, and pepper.
  \item Top the mince mixture with kumara mash. Bake for 25 minutes until golden.
\end{enumerate}

\begin{tipbox}
  \textbf{Tip:} The red lentils bulk out the mince significantly --- you'd never notice they're there, but they add fibre and extra protein cheaply.
\end{tipbox}

\clearpage

% ─── Recipe 3 ────────────────────────────────────────────────────────────────
\recipetitle{3. Beef Mince Stir Fry}
\recipemeta{\$2.50/serve}{470 kcal}{52g}

\ingredientheader
\begin{itemize}[leftmargin=1.5em]
  \item 500g lean beef mince
  \item 1 bag (500g) frozen stir fry vegetables
  \item 200g frozen edamame (shelled)
  \item 3 tbsp soy sauce
  \item 2 tbsp oyster sauce
  \item 1 tbsp sesame oil
  \item 3 cloves garlic, minced
  \item 1 tsp fresh ginger, grated
  \item \textbf{To serve:} jasmine rice or noodles
\end{itemize}

\methodheader
\begin{enumerate}[leftmargin=1.5em]
  \item Cook rice or noodles according to packet instructions.
  \item Heat a wok or large frying pan over high heat. Add a splash of oil.
  \item Cook beef mince, breaking it up, until browned (about 5--7 minutes).
  \item Add garlic and ginger, stir fry for 1 minute.
  \item Add frozen vegetables and edamame and stir fry for 3--4 minutes.
  \item Pour over soy sauce, oyster sauce, and sesame oil. Toss to combine.
  \item Serve immediately over rice or noodles.
\end{enumerate}

\begin{tipbox}
  \textbf{Tip:} Pork mince also works great here for a different flavour. Day-old rice works best for fried rice variants.
\end{tipbox}

\clearpage

% ─── Recipe 4 ────────────────────────────────────────────────────────────────
\recipetitle{4. Slow Cooker Chicken \& Chickpea Stew}
\recipemeta{\$2.50/serve}{500 kcal}{46g}

\ingredientheader
\begin{itemize}[leftmargin=1.5em]
  \item 1 kg chicken thighs
  \item 2 cans chickpeas, drained and rinsed
  \item 1 can diced tomatoes
  \item 1 can coconut milk
  \item 2 tbsp medium curry paste
  \item 2 cups baby spinach
  \item 3 cloves garlic, minced
  \item 1 tsp fresh ginger, grated
  \item \textbf{To serve:} rice or naan bread
\end{itemize}

\methodheader
\begin{enumerate}[leftmargin=1.5em]
  \item Add chicken thighs, chickpeas, diced tomatoes, coconut milk, curry paste, garlic, and ginger to the slow cooker.
  \item Stir to combine. Cook on \textbf{low for 7--8 hours}.
  \item In the last 10 minutes, stir through baby spinach until wilted.
  \item Serve over rice or with warm naan bread.
\end{enumerate}

\begin{tipbox}
  \textbf{Tip:} Use Panang or Tikka Masala curry paste for a medium heat. Rogan Josh will be spicier.
\end{tipbox}

\clearpage

% ─── Recipe 5 ────────────────────────────────────────────────────────────────
\recipetitle{5. Beef \& Vegetable Chilli}
\recipemeta{\$2.60/serve}{490 kcal}{45g}

\ingredientheader
\begin{itemize}[leftmargin=1.5em]
  \item 500g lean beef mince
  \item 1 can kidney beans, drained
  \item 1 can black beans, drained
  \item 2 cans diced tomatoes
  \item 1 can corn, drained
  \item 1 capsicum, diced
  \item 1 onion, diced
  \item 2 tsp chilli powder (medium)
  \item 1 tsp cumin
  \item 1 tsp smoked paprika
  \item 3 cloves garlic, minced
  \item \textbf{To serve:} rice or baked potato, Greek yogurt
\end{itemize}

\methodheader
\begin{enumerate}[leftmargin=1.5em]
  \item Brown mince in a large pot over medium-high heat. Drain excess fat.
  \item Add onion, capsicum, and garlic. Cook for 5 minutes.
  \item Add all remaining ingredients. Stir to combine.
  \item Simmer on low heat for 30--40 minutes, stirring occasionally.
  \item Serve over rice or a baked potato with Greek yogurt on top.
\end{enumerate}

\begin{tipbox}
  \textbf{Tip:} This freezes brilliantly --- make a double batch and freeze half in portions for a quick future dinner.
\end{tipbox}

\clearpage

% ─── Recipe 6 ────────────────────────────────────────────────────────────────
\recipetitle{6. Slow Cooker Pulled Pork Wraps}
\recipemeta{\$2.60/serve}{495 kcal}{54g}

\ingredientheader
\begin{itemize}[leftmargin=1.5em]
  \item 1 kg pork shoulder
  \item 1 tbsp smoked paprika
  \item 1 tsp cumin
  \item 1 tsp garlic powder
  \item 1 tsp brown sugar
  \item 1 cup chicken stock
  \item 1 can black beans, drained and rinsed
  \item \textbf{To serve:} large flour tortillas, coleslaw, Greek yogurt
\end{itemize}

\methodheader
\begin{enumerate}[leftmargin=1.5em]
  \item Mix smoked paprika, cumin, garlic powder, and brown sugar together to form a dry rub.
  \item Rub the spice mixture all over the pork shoulder.
  \item Place in slow cooker and pour chicken stock around the base.
  \item Cook on \textbf{low for 8 hours}.
  \item Shred the pork using two forks. Stir through the cooking juices.
  \item Warm black beans in a small pan with a pinch of cumin and smoked paprika.
  \item Serve in warmed tortillas with black beans, coleslaw and a dollop of Greek yogurt.
\end{enumerate}

\begin{tipbox}
  \textbf{Tip:} Pork shoulder is one of the cheapest cuts available. Buy it on special and freeze --- it slow cooks perfectly from thawed.
\end{tipbox}

\clearpage

% ─── Recipe 7 ────────────────────────────────────────────────────────────────
\recipetitle{7. Slow Cooker Honey Garlic Chicken}
\recipemeta{\$2.40/serve}{480 kcal}{46g}

\ingredientheader
\begin{itemize}[leftmargin=1.5em]
  \item 1 kg chicken thighs
  \item 1 can cannellini beans, drained and rinsed
  \item 4 cloves garlic, minced
  \item 3 tbsp honey
  \item 3 tbsp soy sauce
  \item 1 tbsp apple cider vinegar
  \item 1 tsp sesame oil
  \item 1 tsp cornflour mixed with 2 tbsp water (to thicken)
  \item \textbf{To serve:} steamed rice, steamed broccoli or bok choy
\end{itemize}

\methodheader
\begin{enumerate}[leftmargin=1.5em]
  \item Mix garlic, honey, soy sauce, apple cider vinegar, and sesame oil together in a bowl.
  \item Place chicken thighs and cannellini beans in the slow cooker. Pour sauce over the top.
  \item Cook on \textbf{low for 6--7 hours}.
  \item In the last 20 minutes, stir in the cornflour mixture to thicken the sauce.
  \item Serve over steamed rice with broccoli or bok choy on the side.
\end{enumerate}

\begin{tipbox}
  \textbf{Tip:} The cannellini beans soak up the honey garlic sauce beautifully and add extra protein without changing the flavour at all.
\end{tipbox}

\clearpage

% ─── Recipe 8 ────────────────────────────────────────────────────────────────
\recipetitle{8. Pork \& Lentil Bolognese}
\recipemeta{\$2.30/serve}{460 kcal}{44g}

\ingredientheader
\begin{itemize}[leftmargin=1.5em]
  \item 500g pork mince
  \item 1 cup red lentils, rinsed
  \item 2 cans diced tomatoes
  \item 2 tbsp tomato paste
  \item 1 onion, diced
  \item 2 carrots, diced
  \item 2 stalks celery, diced
  \item 3 cloves garlic, minced
  \item 1 tsp dried oregano
  \item 1 tsp dried basil
  \item 1 cup beef or chicken stock
  \item \textbf{To serve:} pasta or zucchini noodles, parmesan
\end{itemize}

\methodheader
\begin{enumerate}[leftmargin=1.5em]
  \item Brown pork mince in a large pot. Drain excess fat.
  \item Add onion, carrots, celery, and garlic. Cook for 5 minutes until softened.
  \item Stir in tomato paste and cook for 1 minute.
  \item Add diced tomatoes, red lentils, stock, oregano, and basil.
  \item Simmer on low for 30 minutes until lentils are completely soft and sauce has thickened.
  \item Serve over pasta or zucchini noodles with parmesan.
\end{enumerate}

\begin{tipbox}
  \textbf{Tip:} For a lower carb option, zucchini noodles cut the calories significantly while keeping the meal filling.
\end{tipbox}

\clearpage

% ─── Recipe 9 ────────────────────────────────────────────────────────────────
\recipetitle{9. Beef \& Broccoli Stir Fry}
\recipemeta{\$2.70/serve}{465 kcal}{46g}

\ingredientheader
\begin{itemize}[leftmargin=1.5em]
  \item 500g beef strips or beef mince
  \item 1 large head broccoli, cut into florets
  \item 1 can edamame (or 200g frozen shelled edamame)
  \item 1 onion, sliced
  \item 3 cloves garlic, minced
  \item 1 tsp fresh ginger, grated
  \item 3 tbsp soy sauce
  \item 1 tbsp oyster sauce
  \item 1 tbsp sesame oil
  \item 1 tsp cornflour mixed with 2 tbsp water
  \item \textbf{To serve:} jasmine rice
\end{itemize}

\methodheader
\begin{enumerate}[leftmargin=1.5em]
  \item Cook rice according to packet instructions.
  \item Heat a wok over high heat with a splash of oil. Sear beef until browned. Set aside.
  \item In the same wok, stir fry onion, garlic, and ginger for 1--2 minutes.
  \item Add broccoli and edamame. Stir fry for 3--4 minutes until broccoli is just tender.
  \item Return beef to the wok. Add soy sauce, oyster sauce, and sesame oil.
  \item Stir in cornflour mixture and toss everything together until the sauce thickens.
  \item Serve immediately over rice.
\end{enumerate}

\begin{tipbox}
  \textbf{Tip:} Broccoli is one of the most protein-dense vegetables available. Combined with edamame, this stir fry punches well above its weight on protein per dollar.
\end{tipbox}

\clearpage

% ─── Recipe 10 ────────────────────────────────────────────────────────────────
\recipetitle{10. Slow Cooker Beef \& Vegetable Soup}
\recipemeta{\$2.20/serve}{420 kcal}{40g}

\ingredientheader
\begin{itemize}[leftmargin=1.5em]
  \item 500g beef chuck or stewing beef, diced
  \item 1 can cannellini beans, drained and rinsed
  \item 1 can diced tomatoes
  \item 2 carrots, diced
  \item 3 stalks celery, diced
  \item 2 potatoes, diced
  \item 1 onion, diced
  \item 3 cloves garlic, minced
  \item 1 litre beef stock
  \item 1 tsp dried thyme
  \item 1 tsp dried rosemary
  \item 1 tbsp Worcestershire sauce
  \item Large handful of spinach or kale
  \item \textbf{To serve:} crusty bread or dinner rolls
\end{itemize}

\methodheader
\begin{enumerate}[leftmargin=1.5em]
  \item Add all ingredients except spinach to the slow cooker. Stir to combine.
  \item Cook on \textbf{low for 8 hours} or high for 4--5 hours.
  \item In the last 10 minutes, stir through spinach or kale until wilted.
  \item Taste and adjust seasoning. Serve with crusty bread.
\end{enumerate}

\begin{tipbox}
  \textbf{Tip:} This is a great winter warmer and keeps extremely well in the fridge for 4--5 days. The flavour improves the next day. Freeze any extra portions in individual containers.
\end{tipbox}

% ─── Recipe 11 ────────────────────────────────────────────────────────────────
\clearpage
\recipetitle{11. Slow Cooker Pork \& Kumara Curry}
\recipemeta{\$2.50/serve}{490 kcal}{44g}

\ingredientheader
\begin{itemize}[leftmargin=1.5em]
  \item 1 kg pork shoulder, diced
  \item 2 large kumara, peeled and cubed
  \item 1 can coconut milk
  \item 1 can diced tomatoes
  \item 1 can chickpeas, drained and rinsed
  \item 3 tbsp medium curry paste
  \item 1 onion, diced
  \item 3 cloves garlic, minced
  \item 1 tsp fresh ginger, grated
  \item 2 cups baby spinach
  \item \textbf{To serve:} rice or naan bread
\end{itemize}

\methodheader
\begin{enumerate}[leftmargin=1.5em]
  \item Add pork, kumara, coconut milk, diced tomatoes, chickpeas, curry paste, onion, garlic, and ginger to the slow cooker.
  \item Stir to combine. Cook on \textbf{low for 7--8 hours}.
  \item In the last 10 minutes, stir through spinach until wilted.
  \item Serve over rice or with warm naan.
\end{enumerate}

\begin{tipbox}
  \textbf{Tip:} Kumara breaks down slightly during cooking, naturally thickening the curry. If you prefer more texture, add the kumara in the last 2 hours.
\end{tipbox}

\clearpage

% ─── Recipe 12 ────────────────────────────────────────────────────────────────
\recipetitle{12. Beef Mince Stuffed Capsicums}
\recipemeta{\$2.80/serve}{460 kcal}{43g}

\ingredientheader
\begin{itemize}[leftmargin=1.5em]
  \item 6 large capsicums (any colour), tops cut off and deseeded
  \item 500g lean beef mince
  \item 1 cup cooked rice
  \item 1 can black beans, drained and rinsed
  \item 1 can diced tomatoes
  \item 1 onion, diced
  \item 3 cloves garlic, minced
  \item 1 tsp cumin
  \item 1 tsp smoked paprika
  \item 100g cheddar cheese, grated
  \item Salt and pepper
\end{itemize}

\methodheader
\begin{enumerate}[leftmargin=1.5em]
  \item Preheat oven to 190\textdegree C.
  \item Brown mince in a pan over medium-high heat. Drain excess fat.
  \item Add onion and garlic. Cook for 3--4 minutes until softened.
  \item Stir in diced tomatoes, black beans, rice, cumin, paprika, salt, and pepper. Simmer for 5 minutes.
  \item Stand capsicums upright in a baking dish. Fill each one generously with the mince mixture.
  \item Top each capsicum with grated cheese. Add \textonehalf{} cup water to the base of the dish.
  \item Cover with foil and bake for 30 minutes. Remove foil and bake a further 10--15 minutes until capsicums are tender and cheese is golden.
\end{enumerate}

\begin{tipbox}
  \textbf{Tip:} Makes 6 serves --- one stuffed capsicum each. Leftovers reheat well in the microwave for 2 minutes. A great visually impressive meal to cook for your flatmate.
\end{tipbox}

\clearpage

% ─── Recipe 13 ────────────────────────────────────────────────────────────────
\recipetitle{13. Chicken \& Vegetable Bake}
\recipemeta{\$2.60/serve}{455 kcal}{44g}

\ingredientheader
\begin{itemize}[leftmargin=1.5em]
  \item 1 kg chicken thighs
  \item 3 potatoes, cut into wedges
  \item 2 carrots, cut into chunks
  \item 1 zucchini, sliced
  \item 1 capsicum, sliced
  \item 1 can cannellini beans, drained and rinsed
  \item 3 tbsp olive oil
  \item 4 cloves garlic, minced
  \item 1 tsp smoked paprika
  \item 1 tsp dried oregano
  \item 1 tsp dried thyme
  \item Juice of 1 lemon
  \item Salt and pepper
\end{itemize}

\methodheader
\begin{enumerate}[leftmargin=1.5em]
  \item Preheat oven to 200\textdegree C.
  \item Mix olive oil, garlic, paprika, oregano, thyme, lemon juice, salt and pepper in a large bowl.
  \item Add chicken thighs, potato wedges, carrots, and cannellini beans. Toss to coat everything evenly.
  \item Spread across a large roasting dish in a single layer. Roast for 25 minutes.
  \item Add zucchini and capsicum. Return to oven for a further 15--20 minutes until chicken is cooked through and potatoes are golden.
\end{enumerate}

\begin{tipbox}
  \textbf{Tip:} Everything cooks in one pan --- minimal washing up. Prep the vegetables and marinate the chicken the night before to save time on cooking day.
\end{tipbox}

\clearpage

% ─── Recipe 14 ────────────────────────────────────────────────────────────────
\recipetitle{14. Pork Mince Lettuce Cups (San Choy Bow)}
\recipemeta{\$2.30/serve}{410 kcal}{42g}

\ingredientheader
\begin{itemize}[leftmargin=1.5em]
  \item 500g pork mince
  \item 1 can water chestnuts, drained and roughly chopped
  \item 1 can edamame, drained
  \item 2 spring onions, sliced
  \item 3 cloves garlic, minced
  \item 1 tsp fresh ginger, grated
  \item 3 tbsp soy sauce
  \item 1 tbsp oyster sauce
  \item 1 tbsp sesame oil
  \item 1 tsp rice wine vinegar
  \item 1 tsp chilli flakes (optional)
  \item 1 iceberg lettuce, leaves separated
  \item \textbf{To serve:} steamed rice
\end{itemize}

\methodheader
\begin{enumerate}[leftmargin=1.5em]
  \item Heat a wok over high heat with a splash of oil.
  \item Cook pork mince, breaking it up, until browned (5--7 minutes).
  \item Add garlic and ginger. Stir fry for 1 minute.
  \item Add water chestnuts, edamame, soy sauce, oyster sauce, sesame oil, rice wine vinegar, and chilli flakes. Toss together for 2 minutes.
  \item Stir through spring onions.
  \item Spoon filling into lettuce cups and serve immediately with steamed rice on the side.
\end{enumerate}

\begin{tipbox}
  \textbf{Tip:} A fun meal to serve to your flatmate --- the lettuce cups make it feel like a restaurant dish. Water chestnuts add a great crunch and are usually found in the Asian foods aisle at Pak'nSave.
\end{tipbox}

% ════════════════════════════════════════════════════════════════════════════
%  BREAKFAST RECIPES
% ════════════════════════════════════════════════════════════════════════════
\chapter{Breakfast Recipes}

% ─── Breakfast 1 ─────────────────────────────────────────────────────────────
\recipetitle{1. Egg Muffins}
\recipemeta{\$0.80/serve (2 muffins)}{320 kcal}{28g}

\ingredientheader
\begin{itemize}[leftmargin=1.5em]
  \item 10 eggs
  \item 100g diced bacon or ham
  \item 1 capsicum, finely diced
  \item Large handful of spinach, roughly chopped
  \item 100g cherry tomatoes, halved
  \item 50g feta or cheddar, crumbled or grated
  \item Salt, pepper, mixed dried herbs
\end{itemize}

\methodheader
\begin{enumerate}[leftmargin=1.5em]
  \item Preheat oven to 180\textdegree C. Grease a 12-hole muffin tin well.
  \item Whisk eggs together in a large bowl. Season with salt, pepper, and herbs.
  \item Stir in bacon, capsicum, spinach, tomatoes, and cheese.
  \item Pour mixture evenly into the muffin tin, filling each hole about three-quarters full.
  \item Bake for 18--22 minutes until set and lightly golden on top.
  \item Cool completely before storing in an airtight container in the fridge for up to 5 days.
\end{enumerate}

\begin{tipbox}
  \textbf{Reheat:} 60--90 seconds in the microwave. Eat 2 muffins per serving.
\end{tipbox}

\clearpage

% ─── Breakfast 2 ─────────────────────────────────────────────────────────────
\recipetitle{2. Overnight Oats}
\recipemeta{\$1.20/jar}{450 kcal}{35g}

\textit{Makes 5 jars --- prep all at once on Monday evening.}

\ingredientheader

\textit{Per jar:}
\begin{itemize}[leftmargin=1.5em]
  \item \textonehalf{} cup rolled oats
  \item \textonehalf{} cup Greek yogurt (full fat)
  \item 1 cup milk (or almond/oat milk)
  \item 1 scoop vanilla protein powder
  \item 1 tbsp chia seeds
  \item \textbf{Toppings (add in the morning):} banana, mixed berries, honey, granola
\end{itemize}

\methodheader
\begin{enumerate}[leftmargin=1.5em]
  \item Add oats, Greek yogurt, milk, protein powder, and chia seeds to a jar or container. Stir well.
  \item Seal and refrigerate overnight (or at least 4 hours).
  \item In the morning, give it a stir and add your toppings.
  \item Grab and go --- no heating required.
\end{enumerate}

\begin{tipbox}
  \textbf{Tip:} Keeps in the fridge for up to 5 days. Prep all 5 jars on Sunday or Monday to cover the whole week.
\end{tipbox}

\clearpage

% ─── Breakfast 3 ─────────────────────────────────────────────────────────────
\recipetitle{3. Breakfast Burritos}
\recipemeta{\$1.80/burrito}{480 kcal}{38g}

\textit{Makes 6 burritos --- freeze and microwave as needed.}

\ingredientheader
\begin{itemize}[leftmargin=1.5em]
  \item 8 eggs, whisked
  \item 200g chorizo or pork breakfast sausage, sliced
  \item 1 can black beans, drained and rinsed
  \item 1 capsicum, diced
  \item 1 onion, diced
  \item 6 large flour tortillas
  \item 80g cheddar cheese, grated
  \item Salsa and Greek yogurt to serve
  \item Salt, pepper, cumin, smoked paprika
\end{itemize}

\methodheader
\begin{enumerate}[leftmargin=1.5em]
  \item Cook chorizo in a large frying pan over medium heat until browned. Set aside.
  \item In the same pan, saut\'{e} capsicum and onion until softened (about 5 minutes).
  \item Add black beans and season with cumin, paprika, salt, and pepper. Cook for 2 minutes.
  \item Push vegetables to the side. Pour in whisked eggs and scramble until just cooked. Mix everything together.
  \item Lay out tortillas. Divide the filling evenly across all 6. Top with cheese.
  \item Roll tightly into burritos, folding in the sides as you go.
  \item Wrap each burrito individually in foil or cling wrap. Freeze.
\end{enumerate}

\begin{tipbox}
  \textbf{Reheat from frozen:} Remove foil, wrap in a damp paper towel, microwave for 2--3 minutes. Or reheat in foil in oven at 180\textdegree C for 15--20 minutes for a crispier result.
\end{tipbox}

\clearpage

% ─── Breakfast 4 ─────────────────────────────────────────────────────────────
\recipetitle{4. Frittata Slices}
\recipemeta{\$1.10/serve (2 slices)}{380 kcal}{32g}

\textit{Makes 8 slices --- keep in fridge, reheat as needed.}

\ingredientheader
\begin{itemize}[leftmargin=1.5em]
  \item 10 eggs
  \item 150g cooked chicken breast, diced
  \item 1 zucchini, grated (squeeze out excess moisture)
  \item 1 capsicum, diced
  \item Large handful of spinach
  \item 100g feta, crumbled
  \item 2 cloves garlic, minced
  \item 1 tsp dried Italian herbs
  \item Salt and pepper
  \item Olive oil for cooking
\end{itemize}

\methodheader
\begin{enumerate}[leftmargin=1.5em]
  \item Preheat oven to 180\textdegree C.
  \item Heat olive oil in a large oven-safe frying pan over medium heat. Saut\'{e} garlic, capsicum, and zucchini for 4--5 minutes. Add spinach and wilt.
  \item Whisk eggs in a large bowl with Italian herbs, salt, and pepper. Stir in chicken and feta.
  \item Pour egg mixture over the vegetables in the pan. Stir gently to distribute.
  \item Cook on the stovetop for 4--5 minutes until the edges begin to set.
  \item Transfer to the oven and bake for 15--20 minutes until fully set and golden on top.
  \item Cool, then slice into 8 portions. Store in an airtight container in the fridge for up to 5 days.
\end{enumerate}

\begin{tipbox}
  \textbf{Reheat:} 60--90 seconds in the microwave or a few minutes in a dry frying pan. Eat 2 slices per serving.
\end{tipbox}

% ─── Breakfast 5 ─────────────────────────────────────────────────────────────
\clearpage
\recipetitle{5. Banana Protein Pancakes}
\recipemeta{\$1.00/serve (3 pancakes)}{360 kcal}{30g}

\textit{Makes 9 pancakes --- serve 3 per sitting. Batter keeps in fridge 2 days.}

\ingredientheader
\begin{itemize}[leftmargin=1.5em]
  \item 2 ripe bananas, mashed
  \item 4 eggs
  \item 1 scoop vanilla protein powder
  \item \textonehalf{} cup rolled oats (blended into flour)
  \item 1 tsp baking powder
  \item 1 tsp cinnamon
  \item Pinch of salt
  \item \textbf{To serve:} Greek yogurt, sliced banana, honey, mixed berries
\end{itemize}

\methodheader
\begin{enumerate}[leftmargin=1.5em]
  \item Blend oats into a fine flour using a blender or food processor.
  \item Mash bananas well in a large bowl. Whisk in eggs until combined.
  \item Stir in oat flour, protein powder, baking powder, cinnamon, and salt. Mix until smooth.
  \item Heat a non-stick pan over medium heat with a light spray of oil.
  \item Pour small rounds of batter into the pan (about 3 tbsp each). Cook 2--3 minutes per side until golden and cooked through.
  \item Serve with Greek yogurt, berries, and a drizzle of honey.
\end{enumerate}

\begin{tipbox}
  \textbf{Tip:} The riper the bananas, the sweeter the pancakes. Overripe spotty bananas are ideal and cheapest at the supermarket.
\end{tipbox}

\clearpage

% ─── Breakfast 6 ─────────────────────────────────────────────────────────────
\recipetitle{6. Savoury Egg \& Rice Bowl}
\recipemeta{\$1.30/serve}{410 kcal}{34g}

\textit{Great for using up leftover rice from dinner the night before.}

\ingredientheader
\begin{itemize}[leftmargin=1.5em]
  \item 1 cup cooked jasmine rice (day-old works best)
  \item 3 eggs
  \item 100g diced ham or leftover cooked chicken
  \item \textonehalf{} cup frozen peas or edamame
  \item 2 spring onions, sliced
  \item 1 tbsp soy sauce
  \item 1 tsp sesame oil
  \item Sriracha or chilli flakes to taste
\end{itemize}

\methodheader
\begin{enumerate}[leftmargin=1.5em]
  \item Heat a frying pan or wok over medium-high heat with a splash of oil.
  \item Add rice and stir fry for 2--3 minutes until slightly crispy.
  \item Push rice to the side. Crack in 2 eggs and scramble until just set, then mix through the rice.
  \item Add ham or chicken and peas. Toss for 2 minutes.
  \item Season with soy sauce and sesame oil.
  \item Top with a fried egg, spring onions, and sriracha.
\end{enumerate}

\begin{tipbox}
  \textbf{Tip:} This comes together in under 10 minutes and uses up leftovers. A great alternative when you want something warm and savoury in the morning.
\end{tipbox}

\clearpage

% ─── Breakfast 7 ─────────────────────────────────────────────────────────────
\recipetitle{7. High-Protein Chia Pudding}
\recipemeta{\$1.40/jar}{380 kcal}{32g}

\textit{Makes 5 jars --- prep alongside overnight oats on the same evening.}

\ingredientheader

\textit{Per jar:}
\begin{itemize}[leftmargin=1.5em]
  \item 3 tbsp chia seeds
  \item 1 cup milk (or coconut milk for creamier result)
  \item \textonehalf{} cup Greek yogurt
  \item 1 scoop vanilla protein powder
  \item 1 tsp honey
  \item \textbf{Toppings:} mixed berries, sliced mango, toasted coconut, granola
\end{itemize}

\methodheader
\begin{enumerate}[leftmargin=1.5em]
  \item Whisk milk, Greek yogurt, protein powder, and honey together in a bowl until smooth.
  \item Stir in chia seeds. Mix well to prevent clumping.
  \item Pour into jars or containers. Seal and refrigerate for at least 4 hours or overnight.
  \item In the morning, give it a stir --- it should be thick and pudding-like. Add a splash of milk if too thick.
  \item Top with berries, mango, toasted coconut, or granola before eating.
\end{enumerate}

\begin{tipbox}
  \textbf{Tip:} Chia seeds are high in omega-3 fats and fibre, which help keep you full throughout the morning. Available cheaply in bulk at Pak'nSave.
\end{tipbox}

% ════════════════════════════════════════════════════════════════════════════
%  SNACK OPTIONS
% ════════════════════════════════════════════════════════════════════════════
\chapter{Snack Options}

\begin{tipbox}
  \textbf{Why snacks matter:} On two meals a day, snacks are essential for hitting your 150g protein and 2,000+ kcal daily targets. Aim for snacks with at least 15--20g protein each. These options are all quick, cheap, and require minimal or no preparation.
\end{tipbox}

% ─── Snack 1 ─────────────────────────────────────────────────────────────────
\recipetitle{1. Greek Yogurt \& Berry Bowl}
\recipemeta{\$1.20/serve}{220 kcal}{20g}

\ingredientheader
\begin{itemize}[leftmargin=1.5em]
  \item 200g full-fat Greek yogurt
  \item \textonehalf{} cup mixed berries (fresh or frozen/thawed)
  \item 1 tbsp honey
  \item 1 tbsp granola or crushed almonds
\end{itemize}

\methodheader
\begin{enumerate}[leftmargin=1.5em]
  \item Spoon Greek yogurt into a bowl.
  \item Top with berries, a drizzle of honey, and granola or almonds.
  \item Eat immediately.
\end{enumerate}

\begin{tipbox}
  \textbf{Tip:} Buy the 1kg tub of Anchor or Cyclops Greek yogurt at Pak'nSave --- it's significantly cheaper per gram than individual pottles.
\end{tipbox}

\clearpage

% ─── Snack 2 ─────────────────────────────────────────────────────────────────
\recipetitle{2. Boiled Eggs \& Crackers}
\recipemeta{\$0.60/serve}{210 kcal}{18g}

\ingredientheader
\begin{itemize}[leftmargin=1.5em]
  \item 3 eggs
  \item 6 wholegrain crackers
  \item Salt, pepper, and a dash of hot sauce or mustard
\end{itemize}

\methodheader
\begin{enumerate}[leftmargin=1.5em]
  \item Bring a pot of water to the boil. Gently lower in eggs.
  \item Boil for 7--8 minutes for hard-boiled, or 6 minutes for a slightly jammy yolk.
  \item Transfer to cold water for 2 minutes before peeling.
  \item Season with salt, pepper, and hot sauce. Serve with crackers.
\end{enumerate}

\begin{tipbox}
  \textbf{Tip:} Batch boil 6--8 eggs at once and keep them unpeeled in the fridge for up to a week --- perfect grab-and-go snack.
\end{tipbox}

\clearpage

% ─── Snack 3 ─────────────────────────────────────────────────────────────────
\recipetitle{3. Cottage Cheese \& Tomato Toast}
\recipemeta{\$0.90/serve}{240 kcal}{22g}

\ingredientheader
\begin{itemize}[leftmargin=1.5em]
  \item 2 slices wholegrain bread, toasted
  \item \textonehalf{} cup cottage cheese
  \item 1 tomato, sliced
  \item Salt, pepper, dried herbs or chilli flakes
\end{itemize}

\methodheader
\begin{enumerate}[leftmargin=1.5em]
  \item Toast bread until golden.
  \item Spread cottage cheese generously on each slice.
  \item Top with sliced tomato. Season with salt, pepper, and herbs.
\end{enumerate}

\begin{tipbox}
  \textbf{Tip:} Cottage cheese is one of the cheapest high-protein foods available --- around \$4--5 for a 500g tub at Pak'nSave. Also great with cucumber or smashed avocado on top.
\end{tipbox}

\clearpage

% ─── Snack 4 ─────────────────────────────────────────────────────────────────
\recipetitle{4. Tuna \& Crackers}
\recipemeta{\$1.50/serve}{230 kcal}{28g}

\ingredientheader
\begin{itemize}[leftmargin=1.5em]
  \item 1 can tuna in springwater, drained
  \item 8 wholegrain crackers
  \item 1 tbsp Greek yogurt or light mayo
  \item \textonehalf{} tsp lemon juice
  \item Salt, pepper, chilli flakes
  \item Optional: diced cucumber, spring onion
\end{itemize}

\methodheader
\begin{enumerate}[leftmargin=1.5em]
  \item Mix drained tuna with Greek yogurt, lemon juice, salt, pepper, and chilli flakes.
  \item Add diced cucumber or spring onion if using.
  \item Spoon onto crackers and eat immediately.
\end{enumerate}

\begin{tipbox}
  \textbf{Tip:} One of the highest protein-per-dollar snacks available. Buy Sealord or Greenseas tuna in bulk packs at Pak'nSave for the best value.
\end{tipbox}

\clearpage

% ─── Snack 5 ─────────────────────────────────────────────────────────────────
\recipetitle{5. Peanut Butter \& Banana Rice Cakes}
\recipemeta{\$0.80/serve}{250 kcal}{8g}

\ingredientheader
\begin{itemize}[leftmargin=1.5em]
  \item 3 plain rice cakes
  \item 2 tbsp natural peanut butter
  \item 1 banana, sliced
  \item Pinch of cinnamon
\end{itemize}

\methodheader
\begin{enumerate}[leftmargin=1.5em]
  \item Spread peanut butter over each rice cake.
  \item Top with banana slices and a pinch of cinnamon.
\end{enumerate}

\begin{tipbox}
  \textbf{Tip:} Lower protein than other snacks but high in healthy fats and carbs --- great as a pre-walk or post-exercise snack for energy. Add a scoop of protein powder to a glass of milk alongside it to boost protein.
\end{tipbox}

\clearpage

% ─── Snack 6 ─────────────────────────────────────────────────────────────────
\recipetitle{6. Edamame with Sea Salt}
\recipemeta{\$1.00/serve}{190 kcal}{17g}

\ingredientheader
\begin{itemize}[leftmargin=1.5em]
  \item 200g frozen edamame (in pods or shelled)
  \item Sea salt flakes
  \item Optional: soy sauce or chilli flakes
\end{itemize}

\methodheader
\begin{enumerate}[leftmargin=1.5em]
  \item Cook edamame from frozen --- either microwave in a bowl with a splash of water for 3--4 minutes, or boil in salted water for 4--5 minutes.
  \item Drain and sprinkle generously with sea salt.
  \item Serve with a small bowl of soy sauce for dipping if desired.
\end{enumerate}

\begin{tipbox}
  \textbf{Tip:} One of the most satisfying low-calorie, high-protein snacks. Frozen bags are available cheaply at Pak'nSave and keep in the freezer for months.
\end{tipbox}

\clearpage

% ─── Snack 7 ─────────────────────────────────────────────────────────────────
\recipetitle{7. Homemade Protein Balls}
\recipemeta{\$0.50/ball (makes 20)}{130 kcal/ball}{8g/ball}

\ingredientheader
\begin{itemize}[leftmargin=1.5em]
  \item 1 cup rolled oats
  \item \textonehalf{} cup natural peanut butter
  \item 2 scoops vanilla protein powder
  \item 3 tbsp honey
  \item 3 tbsp dark chocolate chips
  \item 2 tbsp chia seeds
  \item 2--4 tbsp milk (to bind if needed)
\end{itemize}

\methodheader
\begin{enumerate}[leftmargin=1.5em]
  \item Mix all ingredients together in a large bowl until well combined.
  \item If the mixture is too dry to hold together, add milk one tablespoon at a time.
  \item Refrigerate the mixture for 20--30 minutes to firm up.
  \item Roll into 20 balls (about the size of a large marble).
  \item Store in an airtight container in the fridge for up to 2 weeks, or freeze for up to 2 months.
\end{enumerate}

\begin{tipbox}
  \textbf{Tip:} Eat 2--3 balls as a snack (260--390 kcal, 16--24g protein). These are great to batch make once a fortnight. Swap peanut butter for almond butter or add desiccated coconut for variety.
\end{tipbox}

\clearpage

% ─── Snack 8 ─────────────────────────────────────────────────────────────────
\recipetitle{8. Hummus \& Vegetable Sticks}
\recipemeta{\$1.10/serve}{200 kcal}{9g}

\ingredientheader
\begin{itemize}[leftmargin=1.5em]
  \item \textonehalf{} cup hummus (store-bought or homemade)
  \item 2 carrots, cut into sticks
  \item 2 stalks celery, cut into sticks
  \item \textonehalf{} capsicum, cut into strips
  \item \textonehalf{} cucumber, cut into sticks
\end{itemize}

\methodheader
\begin{enumerate}[leftmargin=1.5em]
  \item Wash and cut all vegetables into dipping sticks.
  \item Portion hummus into a small container.
  \item Dip and enjoy.
\end{enumerate}

\begin{tipbox}
  \textbf{Tip:} Prep a container of cut vegetables at the start of the week and keep in the fridge --- it removes all friction from choosing this as a snack. Bluebird hummus is cheap and widely available at Pak'nSave.
\end{tipbox}

% ════════════════════════════════════════════════════════════════════════════
%  PROTEIN SHAKE RECIPES
% ════════════════════════════════════════════════════════════════════════════
\chapter{Protein Shake Recipes}

\begin{tipbox}
  \textbf{Equipment:} A blender or NutriBullet gives the best results, but a shaker bottle works for the simpler recipes. All shakes use one scoop of vanilla or unflavoured protein powder (\textasciitilde25g protein per scoop) as a base. Aim for 1--2 shakes per day to bridge your protein gap.
\end{tipbox}

% ─── Shake 1 ─────────────────────────────────────────────────────────────────
\recipetitle{1. Classic Chocolate Peanut Butter Shake}
\recipemeta{\$1.80/shake}{380 kcal}{38g}

\ingredientheader
\begin{itemize}[leftmargin=1.5em]
  \item 1 scoop chocolate or vanilla protein powder
  \item 1 tbsp natural peanut butter
  \item 1 banana (fresh or frozen for extra thickness)
  \item 300ml milk
  \item 1 tsp cocoa powder
  \item 4--5 ice cubes
\end{itemize}

\methodheader
\begin{enumerate}[leftmargin=1.5em]
  \item Add all ingredients to a blender.
  \item Blend on high for 30--45 seconds until smooth and creamy.
  \item Pour into a large glass and drink immediately.
\end{enumerate}

\begin{tipbox}
  \textbf{Tip:} Frozen bananas make the shake thicker and creamier. Peel and freeze overripe bananas in a zip-lock bag --- they last months and are perfect for shakes.
\end{tipbox}

\clearpage

% ─── Shake 2 ─────────────────────────────────────────────────────────────────
\recipetitle{2. Mango \& Coconut Tropical Shake}
\recipemeta{\$1.60/shake}{350 kcal}{32g}

\ingredientheader
\begin{itemize}[leftmargin=1.5em]
  \item 1 scoop vanilla protein powder
  \item 1 cup frozen mango chunks
  \item 200ml coconut milk (light)
  \item 100ml water or milk
  \item Juice of \textonehalf{} lime
  \item 4--5 ice cubes
\end{itemize}

\methodheader
\begin{enumerate}[leftmargin=1.5em]
  \item Add all ingredients to a blender.
  \item Blend until completely smooth.
  \item Pour into a glass and serve immediately.
\end{enumerate}

\begin{tipbox}
  \textbf{Tip:} Frozen mango is much cheaper than fresh and works better in shakes. Available in large bags at Pak'nSave. Swap mango for frozen pineapple for variety.
\end{tipbox}

\clearpage

% ─── Shake 3 ─────────────────────────────────────────────────────────────────
\recipetitle{3. Berry \& Oat Power Shake}
\recipemeta{\$1.50/shake}{400 kcal}{36g}

\ingredientheader
\begin{itemize}[leftmargin=1.5em]
  \item 1 scoop vanilla protein powder
  \item \textonehalf{} cup frozen mixed berries
  \item \textonehalf{} cup rolled oats
  \item 200g Greek yogurt
  \item 250ml milk
  \item 1 tsp honey
  \item 4--5 ice cubes
\end{itemize}

\methodheader
\begin{enumerate}[leftmargin=1.5em]
  \item Add all ingredients to a blender.
  \item Blend on high for 45--60 seconds until completely smooth (oats need a bit longer).
  \item Pour and drink immediately --- this one thickens quickly.
\end{enumerate}

\begin{tipbox}
  \textbf{Tip:} The oats make this shake very filling --- almost meal-replacement level. Great as a post-walk breakfast if you're short on time.
\end{tipbox}

\clearpage

% ─── Shake 4 ─────────────────────────────────────────────────────────────────
\recipetitle{4. Vanilla Cinnamon Shake}
\recipemeta{\$1.20/shake}{310 kcal}{32g}

\ingredientheader
\begin{itemize}[leftmargin=1.5em]
  \item 1 scoop vanilla protein powder
  \item 300ml milk
  \item 1 tsp cinnamon
  \item \textonehalf{} tsp vanilla essence
  \item 1 tsp honey or maple syrup
  \item 4--5 ice cubes
\end{itemize}

\methodheader
\begin{enumerate}[leftmargin=1.5em]
  \item Add all ingredients to a blender or shaker bottle.
  \item Blend or shake vigorously until combined.
  \item Pour over ice and drink immediately.
\end{enumerate}

\begin{tipbox}
  \textbf{Tip:} The simplest shake in the book --- works in a shaker bottle with no blender needed. Great for days when you just want something quick and simple.
\end{tipbox}

\clearpage

% ─── Shake 5 ─────────────────────────────────────────────────────────────────
\recipetitle{5. Green Spinach \& Banana Shake}
\recipemeta{\$1.40/shake}{330 kcal}{30g}

\ingredientheader
\begin{itemize}[leftmargin=1.5em]
  \item 1 scoop vanilla protein powder
  \item 1 large handful of baby spinach
  \item 1 banana (frozen for best texture)
  \item 250ml milk or almond milk
  \item 1 tbsp chia seeds
  \item 1 tsp honey
  \item 4--5 ice cubes
\end{itemize}

\methodheader
\begin{enumerate}[leftmargin=1.5em]
  \item Add spinach and milk to the blender first. Blend briefly to break down the spinach.
  \item Add remaining ingredients and blend until smooth.
  \item Pour and drink immediately.
\end{enumerate}

\begin{tipbox}
  \textbf{Tip:} Don't let the green colour put you off --- you can't taste the spinach at all. It adds a significant micronutrient boost with minimal calories. A great way to sneak in extra vegetables.
\end{tipbox}

% ════════════════════════════════════════════════════════════════════════════
%  HEALTHY DESSERTS \& SWEET TREATS
% ════════════════════════════════════════════════════════════════════════════
\chapter{Healthy Desserts \& Sweet Treats}

\begin{tipbox}
  \textbf{Note:} These treats are designed to satisfy sweet cravings without derailing your progress. They're higher in protein and lower in sugar than typical desserts. Enjoy 1--2 per week as a reward for staying consistent.
\end{tipbox}

% ─── Dessert 1 ───────────────────────────────────────────────────────────────
\recipetitle{1. Frozen Yogurt Bark}
\recipemeta{\$1.00/serve (makes 8 serves)}{120 kcal}{10g}

\ingredientheader
\begin{itemize}[leftmargin=1.5em]
  \item 500g Greek yogurt
  \item 2 tbsp honey
  \item 1 tsp vanilla essence
  \item \textonehalf{} cup mixed berries
  \item 2 tbsp dark chocolate chips
  \item 2 tbsp granola or crushed almonds
\end{itemize}

\methodheader
\begin{enumerate}[leftmargin=1.5em]
  \item Line a baking tray with baking paper.
  \item Mix Greek yogurt, honey, and vanilla together. Spread evenly over the tray about 1cm thick.
  \item Scatter berries, chocolate chips, and granola over the top.
  \item Freeze for at least 3 hours until completely solid.
  \item Break into rough pieces and store in a zip-lock bag in the freezer.
\end{enumerate}

\begin{tipbox}
  \textbf{Tip:} Keeps in the freezer for up to 2 months. A brilliant alternative to ice cream on warm Christchurch evenings.
\end{tipbox}

\clearpage

% ─── Dessert 2 ───────────────────────────────────────────────────────────────
\recipetitle{2. Chocolate Protein Mug Cake}
\recipemeta{\$0.90/serve}{280 kcal}{26g}

\ingredientheader
\begin{itemize}[leftmargin=1.5em]
  \item 1 scoop chocolate protein powder
  \item 2 tbsp cocoa powder
  \item 1 egg
  \item 3 tbsp milk
  \item 1 tbsp honey or maple syrup
  \item \textonehalf{} tsp baking powder
  \item Pinch of salt
  \item 1 tbsp dark chocolate chips (optional)
\end{itemize}

\methodheader
\begin{enumerate}[leftmargin=1.5em]
  \item Mix all dry ingredients together in a large microwave-safe mug.
  \item Add egg, milk, and honey. Mix until smooth with no lumps.
  \item Stir in chocolate chips if using.
  \item Microwave on high for 60--90 seconds. Check at 60 seconds --- it should be just set in the centre.
  \item Allow to cool for 1--2 minutes before eating directly from the mug.
\end{enumerate}

\begin{tipbox}
  \textbf{Tip:} Don't overcook it --- it will go rubbery. It's ready when the top is just set but still slightly moist. Top with a spoonful of Greek yogurt for an extra protein boost.
\end{tipbox}

\clearpage

% ─── Dessert 3 ───────────────────────────────────────────────────────────────
\recipetitle{3. Baked Cinnamon Apples}
\recipemeta{\$0.80/serve}{180 kcal}{5g}

\ingredientheader
\begin{itemize}[leftmargin=1.5em]
  \item 2 large apples, cored
  \item 2 tbsp rolled oats
  \item 1 tbsp honey
  \item 1 tsp cinnamon
  \item 1 tsp butter
  \item Pinch of nutmeg
  \item \textbf{To serve:} Greek yogurt
\end{itemize}

\methodheader
\begin{enumerate}[leftmargin=1.5em]
  \item Preheat oven to 180\textdegree C.
  \item Core apples, leaving the base intact. Place in a small baking dish.
  \item Mix oats, honey, cinnamon, butter, and nutmeg together. Spoon into the hollow of each apple.
  \item Add a splash of water to the base of the dish.
  \item Bake for 25--30 minutes until apples are tender and filling is golden.
  \item Serve warm with a generous dollop of Greek yogurt.
\end{enumerate}

\begin{tipbox}
  \textbf{Tip:} Feels indulgent but is incredibly wholesome. Royal Gala or Granny Smith apples work best. A great dessert for cold evenings.
\end{tipbox}

\clearpage

% ─── Dessert 4 ───────────────────────────────────────────────────────────────
\recipetitle{4. No-Bake Chocolate Oat Cookies}
\recipemeta{\$0.40/cookie (makes 18)}{140 kcal/cookie}{6g/cookie}

\ingredientheader
\begin{itemize}[leftmargin=1.5em]
  \item 2 cups rolled oats
  \item 1 scoop chocolate protein powder
  \item \textonehalf{} cup natural peanut butter
  \item \textonehalf{} cup honey
  \item 3 tbsp cocoa powder
  \item \textonehalf{} cup milk
  \item 1 tsp vanilla essence
  \item Pinch of salt
\end{itemize}

\methodheader
\begin{enumerate}[leftmargin=1.5em]
  \item Combine honey, peanut butter, cocoa powder, and milk in a saucepan over medium heat. Stir until smooth and just beginning to bubble.
  \item Remove from heat immediately. Stir in vanilla and salt.
  \item Add oats and protein powder. Mix until fully combined.
  \item Drop heaped tablespoons onto a baking paper-lined tray.
  \item Refrigerate for 1--2 hours until set firm.
  \item Store in the fridge for up to 2 weeks.
\end{enumerate}

\begin{tipbox}
  \textbf{Tip:} Eat 2--3 cookies as a dessert (280--420 kcal, 12--18g protein). These double as a snack too and are a great alternative to store-bought protein bars which are usually very overpriced.
\end{tipbox}

\clearpage

% ─── Dessert 5 ───────────────────────────────────────────────────────────────
\recipetitle{5. Banana Nice Cream}
\recipemeta{\$0.60/serve}{210 kcal}{8g}

\ingredientheader
\begin{itemize}[leftmargin=1.5em]
  \item 3 ripe bananas, peeled, sliced, and frozen overnight
  \item \textonehalf{} scoop vanilla protein powder
  \item 2 tbsp peanut butter or cocoa powder (choose your flavour)
  \item 2--3 tbsp milk (as needed)
  \item \textbf{Optional toppings:} dark chocolate chips, granola, sliced banana
\end{itemize}

\methodheader
\begin{enumerate}[leftmargin=1.5em]
  \item Add frozen banana slices to a blender or food processor.
  \item Blend, scraping down the sides as needed, until smooth and creamy (2--3 minutes).
  \item Add protein powder and peanut butter or cocoa. Blend again until combined. Add a splash of milk if the mixture is too thick to blend.
  \item Serve immediately as soft-serve, or freeze for 30 minutes for a firmer consistency.
  \item Top with chocolate chips, granola, or fresh banana.
\end{enumerate}

\begin{tipbox}
  \textbf{Tip:} Only works with completely frozen bananas --- the riper and more spotted, the sweeter the result. Keep a bag of frozen bananas in the freezer at all times for both shakes and this dessert.
\end{tipbox}

% ════════════════════════════════════════════════════════════════════════════
%  NUTRITION REFERENCE
% ════════════════════════════════════════════════════════════════════════════
\chapter{Nutrition Reference}

\section*{Daily Targets}
\addcontentsline{toc}{section}{Daily Targets}

\begin{macrobox}
  \begin{tabular}{ll}
    \textbf{Daily Calories:} & 2,000--2,200 kcal \\
    \textbf{Daily Protein:}  & 150g \\
    \textbf{Daily Steps:}    & 8,000 \\
  \end{tabular}\\[8pt]
  Remember to log everything in \textbf{MyFitnessPal} and adjust your portions if your weight loss stalls or you feel overly hungry.
\end{macrobox}

\section*{Recipe Macros at a Glance}
\addcontentsline{toc}{section}{Recipe Macros at a Glance}

\begin{tabularx}{\linewidth}{|X|c|c|c|}
  \hline
  \rowcolor{primary!20}
  \textbf{Recipe} & \textbf{Cost/serve} & \textbf{Calories} & \textbf{Protein} \\
  \hline
  Mexican Chicken          & \$2.20 & 470 kcal & 48g \\
  \hline
  Shepherd's Pie           & \$2.80 & 480 kcal & 42g \\
  \hline
  Beef Mince Stir Fry      & \$2.50 & 470 kcal & 52g \\
  \hline
  Chicken \& Chickpea Stew & \$2.50 & 500 kcal & 46g \\
  \hline
  Beef \& Veg Chilli       & \$2.60 & 490 kcal & 45g \\
  \hline
  Pulled Pork Wraps        & \$2.60 & 495 kcal & 54g \\
  \hline
  Honey Garlic Chicken     & \$2.40 & 480 kcal & 46g \\
  \hline
  Pork \& Lentil Bolo      & \$2.30 & 460 kcal & 44g \\
  \hline
  Beef \& Broccoli Stir Fry & \$2.70 & 465 kcal & 46g \\
  \hline
  Beef \& Veg Soup         & \$2.20 & 420 kcal & 40g \\
  \hline
  Pork \& Kumara Curry     & \$2.50 & 490 kcal & 44g \\
  \hline
  Stuffed Capsicums        & \$2.80 & 460 kcal & 43g \\
  \hline
  Chicken \& Veg Bake      & \$2.60 & 455 kcal & 44g \\
  \hline
  Pork Lettuce Cups        & \$2.30 & 410 kcal & 42g \\
  \hline
  \hline
  \rowcolor{primary!10}
  \multicolumn{4}{|l|}{\textbf{Snacks}} \\
  \hline
  Greek Yogurt Bowl        & \$1.20 & 220 kcal & 20g \\
  \hline
  Boiled Eggs \& Crackers  & \$0.60 & 210 kcal & 18g \\
  \hline
  Cottage Cheese Toast     & \$0.90 & 240 kcal & 22g \\
  \hline
  Tuna \& Crackers         & \$1.50 & 230 kcal & 28g \\
  \hline
  PB \& Banana Rice Cakes  & \$0.80 & 250 kcal & 8g \\
  \hline
  Edamame w/ Sea Salt      & \$1.00 & 190 kcal & 17g \\
  \hline
  Protein Balls (2--3)     & \$1.00--1.50 & 260--390 kcal & 16--24g \\
  \hline
  Hummus \& Veg Sticks     & \$1.10 & 200 kcal & 9g \\
  \hline
  \hline
  \rowcolor{primary!10}
  \multicolumn{4}{|l|}{\textbf{Protein Shakes}} \\
  \hline
  Choc Peanut Butter Shake & \$1.80 & 380 kcal & 38g \\
  \hline
  Mango Coconut Shake      & \$1.60 & 350 kcal & 32g \\
  \hline
  Berry \& Oat Shake       & \$1.50 & 400 kcal & 36g \\
  \hline
  Vanilla Cinnamon Shake   & \$1.20 & 310 kcal & 32g \\
  \hline
  Green Spinach Shake      & \$1.40 & 330 kcal & 30g \\
  \hline
  \hline
  \rowcolor{primary!10}
  \multicolumn{4}{|l|}{\textbf{Desserts \& Sweet Treats}} \\
  \hline
  Frozen Yogurt Bark       & \$1.00 & 120 kcal & 10g \\
  \hline
  Choc Protein Mug Cake    & \$0.90 & 280 kcal & 26g \\
  \hline
  Baked Cinnamon Apples    & \$0.80 & 180 kcal & 5g \\
  \hline
  No-Bake Choc Oat Cookies (2--3) & \$0.80--1.20 & 280--420 kcal & 12--18g \\
  \hline
  Banana Nice Cream        & \$0.60 & 210 kcal & 8g \\
  \hline
  \hline
  \rowcolor{accent!20}
  Egg Muffins (2)          & \$0.80 & 320 kcal & 28g \\
  \hline
  \rowcolor{accent!20}
  Overnight Oats           & \$1.20 & 450 kcal & 35g \\
  \hline
  \rowcolor{accent!20}
  Breakfast Burrito        & \$1.80 & 480 kcal & 38g \\
  \hline
  \rowcolor{accent!20}
  Frittata Slices (2)      & \$1.10 & 380 kcal & 32g \\
  \hline
  \rowcolor{accent!20}
  Banana Protein Pancakes  & \$1.00 & 360 kcal & 30g \\
  \hline
  \rowcolor{accent!20}
  Savoury Egg \& Rice Bowl & \$1.30 & 410 kcal & 34g \\
  \hline
  \rowcolor{accent!20}
  High-Protein Chia Pudding & \$1.40 & 380 kcal & 32g \\
  \hline
\end{tabularx}

\vspace{1cm}

\begin{tipbox}
  \textbf{Hitting 150g protein daily on 2 meals:}\\
  Breakfast (28--38g) + Dinner (42--50g) = roughly 70--88g from meals alone. \\
  Bridge the gap with: 1 protein shake (\textasciitilde25g) + 200g Greek yogurt (\textasciitilde20g) + handful of almonds (\textasciitilde6g) = \textasciitilde51g extra. \\
  Total: \textbf{\textasciitilde140--140g} --- very close to target with minimal extra eating.
\end{tipbox}

% ════════════════════════════════════════════════════════════════════════════
%  SHOPPING TIPS
% ════════════════════════════════════════════════════════════════════════════
\chapter{Shopping Tips}
\addcontentsline{toc}{chapter}{Shopping Tips}

\section*{Budget Shopping in Christchurch}
\addcontentsline{toc}{section}{Budget Shopping in Christchurch}

\begin{itemize}[leftmargin=1.5em, itemsep=6pt]
  \item \textbf{Pak'nSave} for bulk staples --- rice, oats, canned goods, frozen vegetables, and eggs are significantly cheaper here.
  \item \textbf{Chicken thighs over breasts} --- roughly half the price per kilo, more flavour, and perfect for slow cooking.
  \item \textbf{Frozen vegetables} are just as nutritious as fresh and far cheaper. Keep a bag of frozen stir fry mix and frozen spinach always on hand.
  \item \textbf{Buy mince in bulk} when on special. Divide into 500g portions and freeze immediately.
  \item \textbf{Canned legumes} (chickpeas, kidney beans, black beans, lentils) are your best friend --- around \$1.50 a can and dramatically bulk out any meal.
  \item \textbf{Red lentils} are incredibly cheap per kilo and invisible in mince-based dishes --- they add fibre, protein, and volume.
  \item \textbf{Greek yogurt} --- buy the large tub (1kg) rather than individual portions. Cheaper and lasts the week.
  \item \textbf{Eggs} --- buy 18-packs at Pak'nSave for the best value.
\end{itemize}

\section*{Pantry Staples to Always Have}
\addcontentsline{toc}{section}{Pantry Staples to Always Have}

\begin{multicols}{2}
\begin{itemize}[leftmargin=1.5em, itemsep=4pt]
  \item Jasmine rice (2kg bag)
  \item Red lentils (1kg bag)
  \item Rolled oats (1kg)
  \item Canned diced tomatoes
  \item Canned chickpeas, beans \& cannellini beans
  \item Frozen edamame (shelled)
  \item Frozen mixed berries
  \item Frozen mango chunks
  \item Chicken \& beef stock
  \item Soy sauce
  \item Oyster sauce
  \item Sesame oil
  \item Olive oil
  \item Natural peanut butter
  \item Dark chocolate chips
  \item Smoked paprika
  \item Cumin
  \item Garlic powder
  \item Cinnamon
  \item Dried oregano \& basil
  \item Chilli powder
  \item Brown sugar
  \item Worcestershire sauce
  \item Tomato paste
  \item Chia seeds
  \item Large flour tortillas
  \item Wholegrain crackers
  \item Rice cakes
  \item Protein powder (vanilla \& chocolate)
  \item Cocoa powder
  \item Honey
\end{itemize}
\end{multicols}

\vspace{1cm}
\begin{macrobox}
  \textbf{Weekly grocery target: \$60--\$100 NZD}\\[4pt]
  With pantry staples already stocked, your weekly fresh shop should sit comfortably in this range covering all breakfasts and your two cooking nights.
\end{macrobox}

\end{document}
